% !PRINTED 467
\weberpage{467}

%<<CONT
then $A_1=0$, $A_2=0$, $A_3=0$ are the equations of three of these axes. Under the substitution $\chi$ the functions $A_1, A_2, A_3$ pass cyclically into one another. Under the substitution $\Theta$ the functions $A_1, A_2, A_3$ undergo the following permutation:

\[
\begin{pmatrix} A_1, & A_2, & A_3 \\ A_1, & -A_3, & A_2 \end{pmatrix},
\]

as a simple calculation shows, if one actually carries out the substitution $\Theta$ and makes use of the formulae of \S.\,115.

To indicate at least the course of the calculation, let it be remarked that $A_1$ passes, under the substitution $\Theta$ [\S.\,115, (10)], into

\[
\alpha(\gamma^2\varepsilon^2 + \beta\gamma\varepsilon^3 + \beta^2\varepsilon^5)x_1 + \gamma(\alpha^2\varepsilon^6 + \beta^2 + \alpha\beta\varepsilon^2)x_2 + \beta(\alpha\gamma + \gamma^2\varepsilon + \alpha^2\varepsilon^3)x_3
\]

But by \S.\,115, (8), (17), (11):

\[
\begin{aligned}
\gamma^2\varepsilon^2 + \beta\gamma\varepsilon^3 + \beta^2\varepsilon^5 &= \gamma^2\varepsilon^2 + \alpha^2\varepsilon^3 + \beta^2\varepsilon^5 + \alpha\varepsilon^3 \\
&= \alpha(\varepsilon^4+\varepsilon^{-4}) = h(\varepsilon-\varepsilon^{-1}) = \gamma,
\end{aligned}
\]

and hence the coefficient of $x_1$ becomes equal to $\alpha\gamma$, as in $A_1$, and in the same way one transforms the remaining expressions. Now since $\omega$ can be composed from $\Theta$ and $\chi$ (\S.\,72), it follows that the quantities $A_1^2, A_2^2, A_3^2$ are only permuted among themselves by the group $G_{24}$, and that accordingly their symmetric functions are roots of resolvents of degree 7.

The simplest symmetric function of these quantities is the sum

\weq{4}{A_1^2 + A_2^2 + A_3^2,}

and this, multiplied by a suitable numerical factor, we shall introduce as the unknown of the resolvent.

If we order the sum (4) according to $x_1, x_2, x_3$, we obtain for it an expression of the form

\[
\lambda(x_1^2+x_2^2+x_3^2) + \mu(x_2x_3+x_3x_1+x_1x_2),
\]

in which, by (3):

\[
\begin{aligned}
\lambda &= (\alpha^2\gamma^2\varepsilon^{-1} + \beta^2\gamma^2\varepsilon^{-2} + \alpha^2\beta^2\varepsilon^3)\varepsilon \\
\mu &= 2\alpha\beta\gamma(\alpha\varepsilon+\beta\varepsilon^{-3}+\gamma\varepsilon^2)\varepsilon.
\end{aligned}
\]

Here the factor $\varepsilon$ has been taken out of the bracket, in order that the other factor may remain unchanged under the interchange of $\varepsilon$ and $\varepsilon^2$, and may therefore be expressed rationally in terms of $\sqrt{-7}$ (Bd. I, \S.\,171).
